% problem-000205 8 有机化合物基本概念
\section*{8. 有机化合物基本概念}
对题⽬8-1⾄8-5所提问题进⾏判断并提供合理解释，对题⽬8-6和8-7则按所给条件和要求解答。

\subsection*{8-1}
丙酮和六氟丙酮中哪个分⼦偶极矩更⼤？

\subsection*{8-2}
2.4-戊⼆酮的烯醇含量在⽔中还是在正⼰烷中⾼？将溶剂⽔换成⼆甲亚砜，2.4-戊⼆酮的 $\mathsf { p } K _ { a }$ 是降低还是升⾼？

\subsection*{8-3}
$\left( 2 \mathrm { E } , 4 \mathrm { Z } , 6 \mathrm { E } \right)$ -⾟-2,4,6-三烯在光照下关环后两个甲基是顺式的还是反式的？

\subsection*{8-4}
⼄酸⼄酯、⼄酰氯、N,N-⼆甲基⼄酰胺中哪个化合物α-氢酸性最强?

\subsection*{8-5}
吡啶和六氢吡啶哪个分⼦偶极矩更⼤？

\subsection*{8-6}
2-氟-5-三氟甲基吡啶与⼄醇钠/⼄醇的反应速率要⽐ 2.5-⼆氟吡啶快得多。为什么？

\subsection*{8-7}
下表中列出酰胺化反应中常⽤的⼀些活化试剂。画出RCOOH被以下等量试剂活化后的中间体结构式（接下来与胺反应形成酰胺）。

\subsection*{8-7}
-1 PPh3, CBrCl3 8-7-2 $\mathrm { C y N = C = N C y }$ 8-7-3 AcO

（图：）
